% Project status update for supervisor.
% Compile with: xelatex supervisor_update.tex  (Metropolis prefers Xe/LuaLaTeX;
% pdflatex also works but the typography is less crisp). Requires the
% `metropolis` beamer theme.

\documentclass[10pt,aspectratio=169]{beamer}
\usetheme{metropolis}

\usepackage{amsmath, amssymb}
\usepackage{hyperref}

\metroset{progressbar=frametitle, block=fill}

% Math shortcuts
\newcommand{\X}{X}
\newcommand{\A}{A}
\newcommand{\Sx}{S}
\newcommand{\CATE}{\tau}
\newcommand{\cf}{c}
\newcommand{\mzero}{m_{0}}
\newcommand{\SD}{\mathrm{SD}}
\newcommand{\RMST}{\mathrm{RMST}}

\title{FusionForests}
\subtitle{Data fusion for causal survival inference\,---\,project update}
\author{Tijn Jacobs \\ \small Vrije Universiteit Amsterdam}
\date{\today}

\begin{document}

\maketitle

%======================================================================
\begin{frame}{Project and goals}

\textbf{What.}\ A Bayesian nonparametric framework for fusing a
randomised trial (RCT) with real-world data (RWD) under
unmeasured confounding in the RWD. Implemented in the
\texttt{FusionForests} R package.

\medskip
\textbf{Two explicit goals}
\begin{itemize}
  \item[\textbf{(G1)}] \textbf{Efficiency gain} for heterogeneous
        treatment-effect estimation, by borrowing from the RWD.
  \item[\textbf{(G2)}] \textbf{Identifiable extrapolation} of causal
        survival contrasts beyond the RCT follow-up horizon, using
        the longer follow-up of the RWD.
\end{itemize}

\medskip
\textbf{Niche.}\ Existing methods address confounding correction
\emph{or} long-term extrapolation, rarely both, and almost none with a
flexible nonparametric error. The aim is all three at once.

\end{frame}

%======================================================================
\begin{frame}{Model and causal estimands}

\textbf{AFT decomposition} ($Y_i = \log T_i$):
\[
  Y_i \;=\; \mzero(\X_i, \Sx_i)
         \;+\; \A_i\,\CATE(\X_i)
         \;+\; (1-\Sx_i)\,\A_i\,\cf(\X_i)
         \;+\; \sigma\,\varepsilon_i,
\]
with $\mzero, \CATE, \cf$ each modelled by BART forests and
$\varepsilon \sim P$ under a hierarchical-DP prior. $\CATE$ is
identified from the RCT; $\cf$ is identified from the RWD once $\CATE$
is anchored.

\bigskip
\textbf{Three causal contrasts on the survival scale}
\begin{itemize}
  \item $\Delta_{\SD}(t, x)$ --- survival-probability gain at a fixed
        horizon $t$.
  \item $\Delta_{\RMST}(t^{*}, x)$ --- expected life-time gain on
        $[0, t^{*}]$; the headline summary under (G2).
  \item \alert{$\mathrm{AF}(x) = \exp\{\CATE(x)\}$} --- acceleration
        factor: horizon-free, HDP-independent, causally interpretable
        under frailty and effect heterogeneity
        (Brathovde et al.\ 2024).
\end{itemize}

\end{frame}

%======================================================================
\begin{frame}{Progress in the past year}

\textbf{Methodology}
\begin{itemize}
  \item \textbf{HDP nonparametric error}\,---\,five modes from Gaussian
        through to the full hierarchical-DP with per-source recentring.
        The HDP fixes a centring pathology that destabilised the pooled
        DP under bimodal RWD noise.
  \item \textbf{Causal estimands}\,---\,closed forms for
        $\Delta_{\SD}$ and $\Delta_{\RMST}$ under the HDP; full
        argument for the acceleration factor as the principal causal
        target.
  \item \textbf{Posterior projection summaries}
        (Woody--Carvalho--Murray)\,---\,BART posterior of $\CATE$ or
        the AF projected onto a small interpretable basis; uncertainty
        inherited exactly.
  \item \textbf{Framework writeup}\,---\,identification,
        transportability, structural decomposition of the confounding
        function.
\end{itemize}

\textbf{Application}
\begin{itemize}
  \item ACTG175 (RCT) $+$ MACS (RWD): first end-to-end HIV survival
        analysis underway.
\end{itemize}

\textbf{Prior milestone.}\ Horseshoe Forests paper,
\href{https://arxiv.org/abs/2507.22004}{arXiv:2507.22004}.

\end{frame}

%======================================================================
\begin{frame}{Remaining challenges and next steps}

\textbf{Methodological}
\begin{itemize}
  \item \emph{Constancy of $\CATE$ over time}: waning or late-emerging
        effects break extrapolation.
  \item \emph{RCT does double duty}: anchors $\CATE$ \emph{and}
        disambiguates $\CATE$ from $\cf$\,---\,fragile when
        $n_{\mathrm{RCT}}$ is small.
  \item Calendar-time effects in long RWD; informative RWD censoring;
        covariate overlap at long $t$.
\end{itemize}

\textbf{Implementation.}\ HDP combined with per-source scale; adaptive
truncation; DP-aware censored augmentation; cure-fraction;
time-varying $\CATE(x, t)$.

\textbf{Next}
\begin{itemize}
  \item Finish the HIV case study; simulations over
        $n_{\mathrm{RCT}}/n_{\mathrm{RWD}}$ and extrapolation horizons.
  \item Paper draft at \emph{causal fusion} $\cap$
        \emph{long-term extrapolation} $\cap$ \emph{nonparametric
        error}.
\end{itemize}

\end{frame}

\end{document}
